\section{Splitting fields}
\label{section-splitting-fieds}

\noindent
The following lemma is a useful tool for constructing normal field extensions.

\begin{lemma}
\label{lemma-splitting-field}
Let $F$ be a field. Let $P \in F[x]$ be a nonconstant polynomial.
There exists a smallest field extension $E/F$ such that $P$
splits completely over $E$. Moreover, the field extension $E/F$ is normal
and unique up to (nonunique) isomorphism.
\end{lemma}

\begin{proof}
Choose an algebraic closure $\overline{F}$. Then we can write
$P = c (x - \beta_1) \ldots (x - \beta_n)$ in $\overline{F}[x]$, see
Lemma \ref{lemma-algebraically-closed}. Note that $c \in F^*$. Set
$E = F(\beta_1, \ldots, \beta_n)$. Then it is clear that $E$ is
minimal with the requirement that $P$ splits completely over $E$.

\medskip\noindent
Next, let $E'$ be another minimal field extension of $F$ such that
$P$ splits completely over $E'$. Write
$P = c (x - \alpha_1) \ldots (x - \alpha_n)$ with $c \in F$ and
$\alpha_i \in E'$. Again it follows from minimality that
$E' = F(\alpha_1, \ldots, \alpha_n)$. Moreover, if we pick
any $\sigma : E' \to \overline{F}$
(Lemma \ref{lemma-map-into-algebraic-closure})
then we immediately see that $\sigma(\alpha_i) = \beta_{\tau(i)}$
for some permutation $\tau : \{1, \ldots, n\} \to \{1, \ldots, n\}$.
Thus $\sigma(E') = E$. This implies that $E'$ is a normal extension
of $F$ by Lemma \ref{lemma-characterize-normal}
and that $E \cong E'$ as extensions of $F$ thereby finishing the proof.
\end{proof}

\begin{definition}
\label{definition-splitting-field}
Let $F$ be a field. Let $P \in F[x]$ be a nonconstant polynomial.
The field extension $E/F$ constructed in Lemma \ref{lemma-splitting-field}
is called the {\it splitting field of $P$ over $F$}.
\end{definition}

\begin{lemma}
\label{lemma-normal-closure}
\begin{slogan}
Existence of normal closure of finite extensions of fields.
\end{slogan}
Let $E/F$ be a finite extension of fields. There exists a unique
smallest finite extension $K/E$ such that $K$ is normal over $F$.
\end{lemma}

\begin{proof}
Choose generators $\alpha_1, \ldots, \alpha_n$ of $E$ over $F$.
Let $P_1, \ldots, P_n$ be the minimal polynomials of
$\alpha_1, \ldots, \alpha_n$ over $F$. Set $P = P_1 \ldots P_n$.
Observe that $(x - \alpha_1) \ldots (x - \alpha_n)$ divides $P$, since
each $(x - \alpha_i)$ divides $P_i$. Say
$P = (x - \alpha_1) \ldots (x - \alpha_n)Q$.
Let $K/E$ be the splitting field of $P$ over $E$.
We claim that $K$ is the splitting field of $P$ over $F$ as well
(which implies that $K$ is normal over $F$).
This is clear because $K/E$ is generated by the roots of
$Q$ over $E$ and $E$ is generated by the roots of
$(x - \alpha_1) \ldots (x - \alpha_n)$ over $F$, hence
$K$ is generated by the roots of $P$ over $F$.

\medskip\noindent
Uniqueness. Suppose that $K'/E$ is a second smallest extension such that
$K'/F$ is normal. Choose an algebraic closure $\overline{F}$ and an
embedding $\sigma_0 : E \to \overline{F}$. By
Lemma \ref{lemma-map-into-algebraic-closure}
we can extend $\sigma_0$ to $\sigma : K \to \overline{F}$ and
$\sigma' : K' \to \overline{F}$.
By Lemma \ref{lemma-intersect-normal} we see that
$\sigma(K) \cap \sigma'(K')$ is normal over $F$.
By minimality we conclude that $\sigma(K) = \sigma'(K')$.
Thus $\sigma^{-1} \circ \sigma' : K' \to K$ gives an isomorphism
of extensions of $E$.
\end{proof}

\begin{definition}
\label{definition-normal-closure}
Let $E/F$ be a finite extension of fields. The field extension $K/E$
constructed in Lemma \ref{lemma-normal-closure}
is called the {\it normal closure $E$ over $F$}.
\end{definition}

\noindent
One can construct the normal closure inside any given normal extension.

\begin{lemma}
\label{lemma-normal-closure-inside-normal}
Let $L/K$ be an algebraic normal extension.
\begin{enumerate}
\item If $L/M/K$ is a subextension with $M/K$ finite, then there exists
a tower $L/M'/M/K$ with $M'/K$ finite and normal.
\item If $L/M'/M/K$ is a tower with $M/K$ normal and $M'/M$ finite,
then there exists a tower $L/M''/M'/M/K$ with $M''/M$
finite and $M''/K$ normal.
\end{enumerate}
\end{lemma}

\begin{proof}
Proof of (1). Let $M'$ be the smallest subextension of $L/K$ containing $M$
which is normal over $K$. By Lemma \ref{lemma-normal-closure}
this is the normal closure of $M/K$ and is finite over $K$.

\medskip\noindent
Proof of (2). Let $\alpha_1, \ldots, \alpha_n \in M'$ generate $M'$ over $M$.
Let $P_1, \ldots, P_n$ be the minimal polynomials of
$\alpha_1, \ldots, \alpha_n$ over $K$. Let $\alpha_{i, j}$ be the roots
of $P_i$ in $L$. Let $M''  = M(\alpha_{i, j})$. It follows from
Lemma \ref{lemma-normally-generated}
(applied with the set of generators $M \cup \{\alpha_{i, j}\}$)
that $M''$ is normal over $K$.
\end{proof}

\noindent
The following lemma can sometimes be used to prove properties
of the normal closure.

\begin{lemma}
\label{lemma-normal-closure-tensor-product}
Let $L/K$ be a finite extension. Let $M/L$ be the normal
closure of $L$ over $K$. Then there is a surjective map
$$
L \otimes_K L \otimes_K \ldots \otimes_K L \longrightarrow M
$$
of $K$-algebras where the number of tensors can be taken
$[L : K]_s \leq [L : K]$.
\end{lemma}

\begin{proof}
Choose an algebraic closure $\overline{K}$ of $K$.
Set $n = [L : K]_s = |\Mor_K(L, \overline{K})|$ with equality by
Lemma \ref{lemma-separable-degree}.
Say $\Mor_K(L, \overline{K}) = \{\sigma_1, \ldots, \sigma_n\}$.
Let $M' \subset \overline{K}$ be the $K$-subalgebra
generated by $\sigma_i(L)$, $i = 1, \ldots, n$. Then $M'$ is
a field since any $K$-subalgebra of $\overline{K}$ is a field.
Any $K$-algebra map from $M'$ to $\overline{K}$ permutes the $\sigma_i$
so sends $M'$ into and onto $M'$. By construction the field $M'$
is generated by conjugates of elements of $\sigma_1(L)$. Having said
this it follows from Lemma \ref{lemma-characterize-normal}
that $M'$ is normal over $K$ and that it is the
smallest normal subextension of $\overline{K}$ containing
$\sigma_1(L)$. By uniqueness of normal closure we have $M \cong M'$.
Finally, there is a surjective map
$$
L \otimes_K L \otimes_K \ldots \otimes_K L \longrightarrow M',
\quad
\lambda_1 \otimes \ldots \otimes \lambda_n \longmapsto
\sigma_1(\lambda_1) \ldots \sigma_n(\lambda_n)
$$
and note that $n \leq [L : K]$ by definition.
\end{proof}












